% Cap1.tex

\chapter{Introducción}
\label{c1} % la etiqueta para referencias

Desde finales del siglo XVIII, el aumento sostenido de las actividades industriales ha incrementado de manera significativa la concentración de Gases de Efecto Invernadero (GEI) en la atmósfera. Este fenómeno ha intensificado el efecto invernadero natural, alterando progresivamente el equilibrio climático del planeta. Para enfrentar este desafío global, en 1992 se estableció la Convención Marco de las Naciones Unidas sobre el Cambio Climático (CMNUCC), cuyo objetivo es coordinar esfuerzos internacionales para mitigar y adaptarse al cambio climático. Como parte de este marco, se definió la realización anual de las Conferencias de las Partes (COP), instancias de negociación y revisión de compromisos climáticos (\citeA{naciones_unidas_convencion_1992}).\medskip


%protocolo Kioto
La COP realizada el 11 de diciembre de 1998 en Kioto, Japón, generó un gran impacto debido a que fue la primera vez que se propusieron metas cuantificables para reducir los GEI.\footnote{Los gases que fueron parte del acuerdo fueron el dióxido de carbono ($CO_2$), el metano ($CH_4$), el óxido nitroso ($N_2O$), los hidrofluorocarbonos ($HFC$), los perfluorocarbonos ($PFC$) y el hexafluoruro de azufre ($SF_6$).} En esta instancia, se estableció un sistema basado en mecanismos de mercado para regular y tender a disminuir las emisiones de GEI, consistente en la creación de permisos transables, en una cantidad limitada definida por el regulador \footnote{En este trabajo, el término \emph{regulador} se utiliza indistintamente con \emph{agente}, \emph{planificador}, \emph{planificador social} o \emph{ente regulador}, refiriéndose siempre a la autoridad que determina la cantidad total de permisos de emisión del sistema.}, que habilitan a su titular a emitir dioxido de carbono equivalente (CO$_2$e) en su actividad económica, abriendo paso así al sistema denominado \emph{Cap-and-Trade} (\citeA{naciones_unidas_protocolo_1998}). Este sistema apoya la reducción de GEI por medio de la imposición de un límite máximo de emisiones, denominado \textit{Cap}. El \textit{Cap} es fijado en un mercado primario \footnote{Mercado en el cual un planificador social o entidad de gobierno decide y vende un nivel limitado de permisos de emisiones posibles para un período de tiempo}, y posteriormente los distintos consumidores de permisos dentro del sistema pueden transaccionar estos permisos en un mercado secundario (\textit{Trade}). \footnote{Mercado en el cual se transaccionan partes o totalidad de permisos entre los miembros}\footnote{Esto funciona de forma similar al mercado de la bolsa de valores donde las empresas emiten parte de sus empresas al mercado a cierto valor y posteriormente estas se transaccionan}\medskip

La CMNUCC decidida por incrementar los esfuerzos ambientales, el 2015 realiza la COP en París, donde se estableció el objetivo de aumentar la temperatura media mundial muy por debajo de los $2°C$ y mantener el esfuerzo para limitar ese aumento de la temperatura en $1,5°C$ con respecto a los niveles preindustriales (\citeA{naciones_unidas_acuerdo_2015}). Al ser partícipe de esta convención, Chile ha formalizado sus compromisos a través de la Ley Marco del Cambio Climático (LMCC)\footnote{Biblioteca del Congreso Nacional de Chile. Ley Marco del Cambio Climático. Ley N°21455. \url{https://www.bcn.cl/leychile/navegar?idNorma=1177286} (revisada el 12/10/2025).}publicada el año 2022, la cual establece la base jurídica para alcanzar la carbono-neutralidad \footnote{Estado de equilibrio donde se compensan las emisiónes con las absorciones de gases de efecto invernadero.} a más tardar en 2050. La ley define una arquitectura de planificación climática vinculante compuesta por la Estrategia Climática de Largo Plazo, el Presupuesto Nacional de Emisiones, las Contribuciones Determinadas a Nivel Nacional y los planes sectoriales de mitigación y adaptación. Estos instrumentos se rigen por principios como la base científica y la costo-efectividad. Además, incorpora planes regionales, comunales y de recursos hídricos para asegurar coherencia entre la planificación nacional y local.\medskip

Se destaca para efectos de este trabajo, que el artículo 15 de la LMCC, establece la posibilidad de utilizar mecanismos flexibles para el uso de los permisos de emisión, incluyendo el comercio de permisos que establecen las bases para la implementación de un futuro sistema \emph{Cap-and-Trade} en el país. Chile hoy utiliza impuestos para limitar las emisiones de carbono equivalente (CO$_2$e)\footnote{Métrica utilizada para medir y comparar distintos gases de efecto invernadero mediante su equivalencia en dióxido de carbono ($CO_2e$).}, pero el éxito de algunos países utilizando mecanismos como el de \emph{Cap-and-Trade}, hacen relevante la comparación de ambos sistemas. \citeA{anderson_cap-and-trade_2025} explican que en la teoría, cuotas\footnote{Con cuotas se refiere al sistema \emph{Cap-and-Trade}.} e impuestos pueden ser equivalentes para la mitigación de los GEI si se diseñan bien, pero en la práctica no siempre lo son debido a la multiplicidad de equilibrios generando incertidumbre: el impuesto no puede asegurar la cantidad de emisiones en el aire y la cuota no asegura un precio de los permisos a emitir.\medskip

\citeA{anderson_cap-and-trade_2025} mencionan que el sistema de cuotas es más robusto, pues permite alcanzar un objetivo sin necesidad de tanta información del gobierno sobre las empresas o tecnologías. El impuesto, por otro lado, requiere un conocimiento detallado de las funciones de costo y producción de cada agente para fijar correctamente la tasa impositiva. Señalan los autores que, incluso teniendo información completa, puede ser imposible para el gobierno alcanzar sus objetivos ambientales por medio de impuestos.\medskip 

Es importante destacar, sin embargo, que la apertura de los mercados de derechos de emisión requiere de un gran esfuerzo. No solo implica una nueva herramienta de política ambiental, sino también un cambio estructural en el comportamiento de las empresas dentro del equilibrio general de la economía. \citeA{mandel_changes_2009} plantea que la creación de estos mercados modifica las condiciones bajo las cuales las empresas maximizan beneficios y toman decisiones de producción, afectando el equilibrio previo de la economía sin mercados de emisión. Este enfoque resulta relevante para comprender las dinámicas de adaptación de las firmas cuando se introduce un sistema de comercio de emisiones como el \emph{Cap-and-Trade}, y las condiciones necesarias para que la economía alcance nuevamente un equilibrio tras su apertura.\medskip

%bibliografía: \citeA{aldy_eu_2009} 
En la práctica, el inicio de la implementación de un sistema \emph{Cap-and-Trade} en la Unión Europea (UE) enfrentó diversos desafíos, entre ellos, la limitada disponibilidad de información para estimar un Cap adecuado y la carencia de infraestructura institucional en algunos países miembros. Estos problemas fueron abordados progresivamente mediante la recopilación sistemática de datos, el fortalecimiento de las capacidades institucionales y ajustes sucesivos en el diseño del mecanismo, lo que permitió una implementación progresivamente más precisa y coherente del sistema (\citeA{aldy_eu_2009}). Hoy en día, el mecanismo refleja un aprendizaje institucional acumulado: en 2024, la UE logró una reducción emisiones en un $47.6\%$ respecto a los niveles de 2005 (\citeA{international_carbon_action_partnership_eu_2024}), lo que fue posible, en parte, debido a que el mecanismo impulsó una fuerte expansión de energías renovables, particularmente de tecnologías solar y eólica. No obstante, la definición del \textit{Cap} continúa siendo un ejercicio complejo y sujeto a incertidumbre, ya que depende de expectativas sobre la evolución económica, la demanda energética y el cambio tecnológico, así como de consideraciones políticas que condicionan la credibilidad y efectividad ambiental del instrumento (\citeA{aldy_eu_2009}), por lo que para desarrollar el sistema correctamente y obtener resultados deseados, se requiere de un gran compromiso institucional.\medskip

California, por su parte, ha realizado una gestión destacada del sistema \emph{Cap-and-Trade}. Su modelo se originó a partir de la Ley AB 32 de 2006, la cual estableció objetivos de reducción de emisiones de gases de efecto invernadero y encargó a la California Air Resources Board (\textit{CARB}, por sus siglas en ingles) el diseño e implementación de instrumentos regulatorios para alcanzarlos, entre ellos el sistema \emph{Cap-and-Trade}. Este programa comenzó a operar en 2012 y constituye uno de los principales mecanismos de mercado utilizados por el estado para limitar las emisiones\footnote{Programa de límite y comercio de emisiones de California. Disponible en: \url{https://icapcarbonaction.com/en/ets/usa-california-cap-and-trade-program} Consultado el 7/03/2026.}. En la implementación del sistema \emph{Cap-and-Trade}, se destaca la utilización de los ingresos obtenidos por la venta de permisos. Estos son invertidos por el Estado en la comunidad, con un énfasis particular en los sectores más vulnerables. En concreto, una parte significativa de estos recursos se destina a facilitar la adopción de vehículos eléctricos y a mejorar la eficiencia energética de los hogares, contribuyendo directamente a una transición hacia estilos de vida más sostenibles y apoyando la reducción de emisiones de gases de efecto invernadero \footnote{Fundación Nueva Generación (Argentina). \emph{El éxito del cap-and-trade en California: oportunidades para adaptar el modelo en Santa Fe}. Disponible en: \url{https://www.fnga.org.ar/el-\%C3\%A9xito-del-cap-and-trade-en-california-oportunidades-para-adaptar-el-modelo-en-santa-fe}. Consultado el 22/12/2025.}. Este tipo de casos evidencia la importancia de una buena estructura y compromiso institucional para enfrentar de la mejor forma los compromisos ambientales. \medskip

Un caso más cercano y reciente es el de Brasil, que en 2024 aprobó el Sistema Brasileño de Comercio de Emisiones (SBCE). Este marco legal busca establecer un mercado regulado nacional que conviva con mecanismos voluntarios y, eventualmente, se integre con mercados internacionales bajo estándares de integridad y trazabilidad. El SBCE representa un avance significativo en la gobernanza regional y en la evaluación comparada de diseños regulatorios (\citeA{international_carbon_action_partnership_icap_brazilian_2022}).\medskip

Chile, a la fecha, utiliza un impuesto de emisiones de 5 dólares por tonelada de carbono emitida \footnote{Biblioteca del Congreso Nacional de Chile. Reforma Tributaria que Modifica el Sistema de Tributación de la Renta e Introduce Diversos Ajustes en el Sistema Tributario. Ley N°20780.\url{https://www.bcn.cl/leychile/navegar?idNorma=1067194} (Revisada el 17/11/2025).}, cifra que según \citeA{amigo_two_2021}\footnote{Modelo también denominado como "modelo original".} no es suficiente para alcanzar la neutralidad al año 2050, mencionando que una cifra más adecuada sería de 24 dólares. Sin embargo, incluso con una cifra más alta, según lo indicado anteriormente, los impuestos no aseguran cumplir las metas. Es por esto que para una futura implementación de mejores sistemas, el trabajo de \citeA{amigo_two_2021} propone implementar un mecanismo de \emph{Cap-and-Trade} para la industria eléctrica de Chile. Estos autores plantean un modelo de optimización en equilibrio estocástico, donde un planificador social emite permisos de emisión y posteriormente los productores de electricidad participan en un mercado donde compran y venden permisos de emisión.\medskip

El modelo de los autores recién mencionados representa adecuadamente el mercado de los productores en la optimización de sus beneficios, producción y equilibrio de mercado. Sin embargo, el planificador social, podría ser mejorado ya que no logra representar el comportamiento de un regulador que, en la práctica, actúa de forma racional y posee la facultad de optimizar la asignación de recursos (tiempo, estudios o presupuesto) para mejorar la precisión de su estimación. El modelo más bien asume información perfecta para la estimación del presupuesto y un costo en su utilidad que no tiene efectos reales en la resolución del modelo. El estudio realizado por \citeA{munoz_lhuillier_inversion_2022} busca atender esta limitación, por lo que se contemplará su trabajo a lo largo del presente texto con fines comparativos.\medskip

\section{Objetivos}
Para desarrollar y evaluar una reformulación del planificador social en el modelo \textit{Cap-and-Trade} de \citeA{amigo_two_2021} se establecen los siguientes objetivos, con el fin de enfocar el estudio.\medskip

\subsection{Objetivo general}
El objetivo general es el de incorporar atención en la conducta del planificador social en el marco de un sistema \textit{Cap-and-Trade} en Chile, evaluando cómo distintos niveles de atención afectan la transición hacia energías limpias y la precisión del presupuesto de emisiones frente a las metas nacionales. \medskip
\subsection{Objetivos Específicos}
Para cumplir este objetivo se establecen objetivos específicos que sirven como guía para el desarrollo del objetivo general:
\begin{enumerate}
    \item Analizar críticamente el rol del subastador en el modelo original y el realizado por \citeA{munoz_lhuillier_inversion_2022} con el objetivo de identificar sus limitaciones teóricas y prácticas.
    \item Diseñar una nueva formulación del subastador incorporando inatención conductual.
    \item Obtener los límites para el $Cap$ coherentes con valores históricos de Chile, por medio de investigación. 
    \item Comparar cuantitativamente los resultados del modelo reformulado con los obtenidos por \citeA{amigo_two_2021} y \citeA{munoz_lhuillier_inversion_2022}.
\end{enumerate}
\medskip
\section{Alcances}
Para asegurar la validez del estudio, se delimitan las fronteras de la investigación. La cual se centra en el límite del presupuesto de carbono definido por Chile para el periodo 2020-2030, debido a que representa el ciclo de cumplimiento de metas climáticas más exigente a la fecha. El estudio se enfoca en la reformulación del planificador social propuesto por \citeA{amigo_two_2021}, manteniendo inalterada la representación del mercado de productores. Esto permite aislar y cuantificar el efecto directo de la inatención conductual en las decisiones regulatorias del sistema original. Para la implementación, el modelo se ejecuta en un entorno de programación desarrollado por \citeA{munoz_lhuillier_inversion_2022} en lenguaje \textit{Julia}, utilizando el \textit{solver} \textit{PATH} debido a su alta eficiencia en la resolución de problemas de complementariedad mixta (MCP).

\medskip

\section{Metodología}
La metodología seguida en este trabajo se compone de las siguientes etapas:

\begin{enumerate}
    \item {Identificación de funciones de costos}: se estudian distintas funciones de costos presentes en la literatura, seleccionando aquella que resulte aplicable y coherente con los objetivos del presente trabajo.
    \item {Formulación matemática del planificador social}: se diseña la nueva formulación incorporando la función de costos seleccionada. Esta se implementa.
    \item {Revisión bibliográfica}: se realiza una búsqueda y análisis de trabajos recientes con el fin de encontrar los parámetros a integrar en la nueva formulación del planificador. 
    \item {Implementación}: se computa la nueva formulación del planificador en un código lenguaje \textit{Julia} previamente desarrollado por \citeA{munoz_lhuillier_inversion_2022}.
    \item {Calibración}: se ajustan los parámetros del modelo para que la matriz energética simulada represente lo más posible a la realidad del sector eléctrico chileno.
    \item {Evaluación del modelo}: se comparan los resultados obtenidos con los modelos de referencia (\citeA{amigo_two_2021} y el de \citeA{munoz_lhuillier_inversion_2022}), con el objetivo de evaluar mejoras en precisión y representatividad.
\end{enumerate}
\medskip
\section{Estructura del documento}

El presente documento presenta cinco Capítulos: Introducción, Marco Teórico, Desarrollo Metodológico, Análisis de Resultados y Conclusiones. En la introducción se da contexto al cambio climático y la necesidad de mejorar los sistemas para reducir las emisiones de GEI. En el marco teórico, se explica el modelo de \citeA{amigo_two_2021} y el de \citeA{munoz_lhuillier_inversion_2022}, donde se muestra a través de gráficos explicativos, la necesidad de reformular el problema del planificador social. Se propone una posible implementación de mejora a los modelos mencionados a través de teoría de atención racional expuesta por \citeA{gabaix_behavioral_2019}. En el desarrollo metodológico se da contexto a la metodología para estimar presupuestos de carbono, se muestran datos históricos de presupuestos para Chile. Se intenta encontrar, además, los parámetros que alimentan el modelo, metodología de calibración y como se implementará el modelo en código de lenguaje \textit{Julia}. En el análisis de los resultados se muestran los datos que entrega el nuevo modelo implementado y análisis de estos. Finalmente, en la conclusión se analizan los resultados, se comparan con los modelos de \citeA{amigo_two_2021} y \citeA{munoz_lhuillier_inversion_2022}, y se propone un espacio para próximos estudios.


